% ============================================================
\chapter{Bonnet-Myers and Synge Theorems}
\label{ch:bonnet_myers}
% ============================================================

\section{Introduction and motivation}

Can one deduce the global shape of a universe from local measurements of its curvature? This is precisely what the Bonnet--Myers and Synge theorems accomplish, two jewels of global Riemannian geometry. Bonnet's theorem (1855), rediscovered and generalized by Myers in 1941, asserts that a complete manifold whose Ricci curvature is strictly positive is necessarily compact and of finite diameter. In other words, positive curvature ``closes'' the space upon itself. Synge's theorem (1936) goes further: a compact Riemannian manifold of even dimension with strictly positive sectional curvature is simply connected. These results illustrate spectacularly how curvature controls topology.

\begin{intuition}
The guiding principle is that positive curvature ``brings geodesics closer
together,'' preventing the manifold from extending to infinity. More precisely:
\begin{itemize}
    \item Ricci curvature bounded below by a positive constant forces a
          \textbf{finite diameter}, hence compactness.
    \item Strictly positive sectional curvature in even dimensions constrains
          \textbf{orientability} and \textbf{simple connectivity}.
\end{itemize}
\end{intuition}

\section{Jacobi fields and Ricci curvature}

\begin{definition}[Ricci curvature]
\index{Ricci curvature}
Let $(M, g)$ be a Riemannian manifold of dimension $n$. The \emph{Ricci curvature}
is the trace of the curvature tensor:
\[
    \operatorname{Ric}(u, v) = \sum_{i=1}^{n} g(R(e_i, u)v, e_i),
\]
where $(e_1, \ldots, e_n)$ is an orthonormal basis of $T_pM$.
\end{definition}

\begin{remark}
The Ricci curvature $\operatorname{Ric}(v, v)$ for $\norm{v} = 1$ is the sum of
sectional curvatures $K(\sigma_i)$ of the planes spanned by $v$ and the remaining
basis vectors. It is thus an ``average'' of sectional curvature.
\end{remark}

\begin{definition}[Jacobi field]
\index{Jacobi field}
Let $\gamma : [0, L] \to M$ be a geodesic. A \emph{Jacobi field} along $\gamma$
is a vector field $J$ along $\gamma$ satisfying the Jacobi equation:
\[
    \frac{D^2 J}{dt^2} + R(J, \dot{\gamma})\dot{\gamma} = 0.
\]
\end{definition}

\begin{definition}[Index form]
\label{def:index_form}
The \emph{index form} of a geodesic $\gamma : [0, L] \to M$ is:
\[
    I(V, W) = \int_0^L \left[ g\!\left(\frac{DV}{dt}, \frac{DW}{dt}\right)
    - g(R(V, \dot{\gamma})\dot{\gamma}, W) \right] dt,
\]
defined on vector fields along $\gamma$ vanishing at the endpoints.
\end{definition}

\section{The Bonnet-Myers theorem}

\begin{theorem}[Bonnet-Myers]
\label{thm:bonnet_myers}
\index{Bonnet-Myers theorem}
Let $(M, g)$ be a complete Riemannian manifold of dimension $n \geq 2$. If the
Ricci curvature satisfies:
\[
    \operatorname{Ric}(v, v) \geq (n-1)\kappa > 0
    \quad \forall\, v \in TM,\; \norm{v} = 1,
\]
then:
\begin{enumerate}
    \item The diameter satisfies $\operatorname{diam}(M) \leq \frac{\pi}{\sqrt{\kappa}}$.
    \item $M$ is compact.
    \item The fundamental group $\pi_1(M)$ is finite.
\end{enumerate}
\end{theorem}

\begin{proof}
\textbf{Step 1: Diameter bound.}
Let $p, q \in M$ and $\gamma : [0, L] \to M$ be a minimizing geodesic from $p$ to
$q$, parametrized by arc length. Suppose for contradiction that
$L > \pi / \sqrt{\kappa}$.

Choose an orthonormal basis $(e_1, \ldots, e_{n-1}, \dot{\gamma}(0))$ of $T_pM$
and let $E_i(t) = P_{0,t}(e_i)$ be the parallel translates along $\gamma$.
Define the variation fields:
\[
    V_i(t) = \sin\!\left(\frac{\pi t}{L}\right) E_i(t).
\]
Each $V_i$ vanishes at $t = 0$ and $t = L$. The index form gives:
\[
    I(V_i, V_i) = \int_0^L \left[ \frac{\pi^2}{L^2} \cos^2\!\left(\frac{\pi t}{L}\right)
    - \sin^2\!\left(\frac{\pi t}{L}\right) K(E_i, \dot{\gamma}) \right] dt.
\]
Summing over $i = 1, \ldots, n-1$:
\[
    \sum_{i=1}^{n-1} I(V_i, V_i) = \int_0^L \sin^2\!\left(\frac{\pi t}{L}\right)
    \left[ \frac{(n-1)\pi^2}{L^2} - \operatorname{Ric}(\dot{\gamma}, \dot{\gamma}) \right] dt.
\]
If $\operatorname{Ric}(\dot{\gamma}, \dot{\gamma}) \geq (n-1)\kappa$ and
$L > \pi/\sqrt{\kappa}$, then $(n-1)\pi^2/L^2 < (n-1)\kappa$, making the sum
strictly negative. Hence at least one $I(V_i, V_i) < 0$, meaning $\gamma$ is not
minimizing --- contradiction.

\textbf{Step 2: Compactness.}
Since $M$ is complete with finite diameter, it is bounded. By the Hopf-Rinow theorem,
closed bounded subsets of a complete manifold are compact, so $M$ is compact.

\textbf{Step 3: Finiteness of $\pi_1(M)$.}
The universal cover $\widetilde{M}$ inherits the pullback metric, which satisfies
the same Ricci curvature bound. By Steps 1--2, $\widetilde{M}$ is compact.
The covering $\pi : \widetilde{M} \to M$ is therefore finite, so
$\abs{\pi_1(M)} < \infty$.
\end{proof}

\begin{warning}[title=\textbf{Completeness is essential}]
Completeness is crucial: an open disk with a positively curved metric is not
compact. Likewise, the \emph{strict} positivity $\kappa > 0$ is essential:
a paraboloid has $\operatorname{Ric} \geq 0$ but is non-compact.
\end{warning}

\begin{corollary}[Comparison with the sphere]
If $(M^n, g)$ is complete with $\operatorname{Ric} \geq (n-1)$, then
$\operatorname{diam}(M) \leq \pi = \operatorname{diam}(S^n)$.
\end{corollary}

\begin{example}[Sphere and projective spaces]
The sphere $S^n$ of radius $r$ has $\operatorname{Ric} = \frac{n-1}{r^2} g$,
so $\kappa = 1/r^2$ and Bonnet-Myers gives $\operatorname{diam}(S^n(r)) \leq \pi r$.
Equality holds: the diameter of $S^n(r)$ is exactly $\pi r$.
\end{example}

\begin{example}[Compact Lie groups]
Every compact semisimple Lie group with the negative of the Killing form has
$\operatorname{Ric} > 0$. Thus $SU(n)$, $SO(n)$, $Sp(n)$ are compact with
finite fundamental group.
\end{example}

\section{Cheng's maximal diameter theorem}

\begin{theorem}[Cheng]
\label{thm:cheng}
\index{Cheng's theorem}
Let $(M^n, g)$ be a complete Riemannian manifold with
$\operatorname{Ric} \geq (n-1)\kappa > 0$. If $\operatorname{diam}(M) = \pi/\sqrt{\kappa}$,
then $M$ is isometric to the sphere $S^n(1/\sqrt{\kappa})$.
\end{theorem}

\begin{remark}
This rigidity theorem shows that the Bonnet-Myers bound is sharp and only the
sphere achieves it.
\end{remark}

\section{Synge's theorem}

\begin{theorem}[Synge]
\label{thm:synge}
\index{Synge's theorem}
Let $(M, g)$ be a compact Riemannian manifold with strictly positive sectional
curvature ($K > 0$). Then:
\begin{enumerate}
    \item If $\dim M$ is even and $M$ is orientable, then $M$ is simply connected.
    \item If $\dim M$ is odd, then $M$ is orientable.
\end{enumerate}
\end{theorem}

\begin{proof}[Proof (even-dimensional case)]
Suppose $\dim M = 2m$ and $M$ is orientable. Assume for contradiction that
$\pi_1(M) \neq \{e\}$. Then there exists a closed geodesic
$\gamma : [0, L] \to M$ of minimal length in its nontrivial free homotopy class.

Parallel transport along $\gamma$ defines an isometry
$P : T_{\gamma(0)}M \to T_{\gamma(0)}M$ preserving $\dot{\gamma}(0)$.
The restriction $P|_{\dot{\gamma}(0)^\perp}$ is an isometry of the
$(2m-1)$-dimensional (odd-dimensional) subspace. Since $M$ is orientable,
$\det(P) = 1$, and the restriction also has determinant $1$. An isometry
of an odd-dimensional space with determinant $1$ has a fixed unit vector $v$.

The parallel vector field $V(t)$ along $\gamma$ with $V(0) = v$ satisfies
$V(L) = P(v) = v = V(0)$, so $V$ is a periodic parallel field perpendicular
to $\dot{\gamma}$. Computing the index form:
\[
    I(V, V) = -\int_0^L K(V, \dot{\gamma}) \, dt < 0,
\]
since $K > 0$. This means $\gamma$ can be shortened by a variation in the
direction of $V$, contradicting the minimality of $\gamma$ in its homotopy class.
\end{proof}

\begin{corollary}
The only even-dimensional manifolds with strictly positive sectional curvature
and nontrivial fundamental group are the real projective spaces $\R P^{2m}$
(which are non-orientable).
\end{corollary}

\begin{example}[$S^{2m}$ and $\C P^m$]
The sphere $S^{2m}$ ($K = 1$) is simply connected, in agreement with Synge.
The complex projective space $\C P^m$ ($K \in [1, 4]$) is orientable of even
dimension $2m$, and indeed $\pi_1(\C P^m) = 0$.
\end{example}

\begin{example}[$S^3$ and $\R P^3$]
In odd dimensions, Synge does not force simple connectivity: $\R P^3$ has $K > 0$
and $\pi_1(\R P^3) = \Z/2\Z \neq 0$, but it is orientable, consistent with the theorem.
\end{example}

\section{Applications and topological consequences}

\begin{proposition}[Constraints on the fundamental group]
Let $(M^n, g)$ be complete with $\operatorname{Ric} \geq (n-1)\kappa > 0$. Then:
\begin{enumerate}
    \item $\pi_1(M)$ is finite (Bonnet-Myers).
    \item If $n$ is even and $M$ is orientable, $\pi_1(M) = 0$ (Synge).
    \item If $n$ is odd and $K > 0$, $M$ is orientable (Synge).
\end{enumerate}
\end{proposition}

\begin{theorem}[Bishop-Gromov volume comparison]
If $(M^n, g)$ is complete with $\operatorname{Ric} \geq (n-1)\kappa$, then for all
$p \in M$ and $0 < r \leq R$:
\[
    \frac{\operatorname{Vol}(B(p, R))}{\operatorname{Vol}(B(p, r))}
    \leq \frac{V_\kappa(R)}{V_\kappa(r)},
\]
where $V_\kappa(r)$ is the volume of the ball of radius $r$ in the space form of
constant curvature $\kappa$.
\end{theorem}

\begin{keyformulas}[title=\textbf{Summary of curvature and topology}]
\begin{center}
\renewcommand{\arraystretch}{1.3}
\begin{tabular}{lll}
\textbf{Hypothesis} & \textbf{Conclusion} & \textbf{Theorem} \\
\hline
$\operatorname{Ric} \geq (n-1)\kappa > 0$, complete
    & $\operatorname{diam} \leq \pi/\sqrt{\kappa}$, compact
    & Bonnet-Myers \\
$\operatorname{Ric} \geq (n-1)\kappa > 0$, complete
    & $\pi_1(M)$ finite
    & Bonnet-Myers \\
$K > 0$, compact, even dim, orientable
    & $\pi_1(M) = 0$
    & Synge \\
$K > 0$, compact, odd dim
    & $M$ orientable
    & Synge \\
$\operatorname{Ric} \geq (n-1)\kappa$, complete
    & Volume comparison
    & Bishop-Gromov \\
\end{tabular}
\end{center}
\end{keyformulas}

\section{Exercises}

\begin{exercise}
Show that $\R P^n$ with the quotient metric from $S^n$ has $K = 1$. For even $n$,
verify that $\R P^n$ is non-orientable with $\pi_1(\R P^n) = \Z/2\Z$, consistent
with Synge's theorem.
\end{exercise}

\begin{exercise}
Let $(M^n, g)$ be complete with $\operatorname{Ric} \geq (n-1)$. Show that
$\operatorname{Vol}(M) \leq \operatorname{Vol}(S^n)$.
\end{exercise}

\begin{exercise}
Let $M = S^2 \times S^2$ with the product metric. Compute the Ricci curvature and
verify that $\operatorname{Ric} > 0$. Deduce that $S^2 \times S^2$ is compact and
simply connected (by Synge, $\dim = 4$ is even).
\end{exercise}

\begin{exercise}[Optimality of Bonnet-Myers]
Construct a family of metrics on $S^n$ with $\operatorname{Ric} \geq (n-1)\kappa$
for a given $\kappa$, whose diameter approaches $\pi/\sqrt{\kappa}$.
\end{exercise}

\begin{exercise}
Show that the torus $\mathbb{T}^n$ admits no metric with strictly positive Ricci
curvature. \emph{Hint:} use $\pi_1(\mathbb{T}^n) = \Z^n$.
\end{exercise}

\begin{exercise}
Let $(M, g)$ be a compact $3$-manifold with $K > 0$. Show that $M$ is orientable.
Can one conclude that $M$ is simply connected? Justify using the example of $\R P^3$.
\end{exercise}
